%%==============================================================================
%%  EmoLearn — IEEE conference paper (IEEEtran), English
%%  Aligned with project state: FICHE_RESUME_ENCADRANTE.md (2026)
%%  Compile: pdflatex → bibtex → pdflatex → pdflatex
%%==============================================================================
\documentclass[conference]{IEEEtran}

\usepackage[utf8]{inputenc}
\usepackage[T1]{fontenc}
\usepackage[english]{babel}
\usepackage[numbers,sort&compress]{natbib}
\usepackage{amsmath,amssymb,amsfonts}
\usepackage{graphicx}
\usepackage{textcomp}
\usepackage{xcolor}
\usepackage{booktabs}
\usepackage{multirow}
\usepackage{array}
\usepackage{url}
\usepackage{balance}
\usepackage{enumitem}
\usepackage[hidelinks,breaklinks=true,bookmarks=true,
            pdfstartview=FitH,
            pdfauthor={EmoLearn},
            pdftitle={EmoLearn: Multimodal Adaptive Tutoring}]{hyperref}

\setlist[itemize]{leftmargin=*,itemsep=1pt,topsep=2pt}
\setlist[enumerate]{leftmargin=*,itemsep=1pt,topsep=2pt}

\def\BibTeX{{\rm B\kern-.05em{\sc i\kern-.025em b}\kern-.08em
    T\kern-.1667em\lower.7ex\hbox{E}\kern-.125emX}}

\begin{document}

\title{EmoLearn: Multimodal Adaptive Tutoring with\\
FER+AU, WavLM Speech, and LLM Content Generation}

\author{
\IEEEauthorblockN{Firstname Lastname}
\IEEEauthorblockA{
\textit{University / Institution}\\
Department of Computer Science\\
City, Country\\
email@example.com}
\and
\IEEEauthorblockN{Firstname Lastname (Supervisor)}
\IEEEauthorblockA{
\textit{University / Institution}\\
Department of Computer Science\\
City, Country\\
supervisor@example.com}
}

\maketitle

%------------------------------------------------------------------------------
\begin{abstract}
Online learning platforms rarely sense when a learner becomes bored, confused,
or frustrated. This paper presents \textit{EmoLearn}, a master-thesis research
prototype that estimates engagement from \textbf{face} and \textbf{voice},
decides whether to intervene, and generates adaptive pedagogical content with a
large language model (LLM) grounded by retrieval-augmented generation (RAG).
In the current product runtime, facial state is decided by a \textbf{FER+AU}
pipeline (pretrained Vision Transformer emotion scores concatenated with ten
MediaPipe action-unit proxies, classified by a Random Forest), while speech uses
a \textbf{WavLM} head fine-tuned to four \textsc{DAiSEE} states with an
IEMOCAP emotion-recognition hybrid anchor. Voice clips for the product path are
selected by DAiSEE-state filtering plus cosine similarity on a 12-dimensional
acoustic/affective profile. Fusion applies a $60\%/40\%$ face/voice prior
reweighted by confidence. Beyond single-clip analysis, the system supports
multi-clip sequences with temporal rules and an affect-aware conversational
tutor. We report component metrics, an aligned audiovisual fusion testbed
($n{=}80$), product-proxy evaluation ($n{=}50$), acoustic pairing A/B results,
and automatic form/similarity measures on generated content, together with an
explicit account of simulation limits and pending expert validation.
\end{abstract}

\begin{IEEEkeywords}
adaptive learning, learner engagement, multimodal fusion, facial action units,
speech emotion recognition, WavLM, large language models, RAG, DAiSEE, IEMOCAP
\end{IEEEkeywords}

%==============================================================================
\section{Introduction}
%==============================================================================

Massive Open Online Courses (MOOCs) and large-scale e-learning platforms have
democratized access to education, yet they continue to exhibit high dropout
rates~\cite{breslow2013studying,daniel2012making,onah2014dropout,jordan2015massive}.
A central driver is declining \textbf{learner
engagement}~\cite{kizilcec2013deconstructing,hew2016promoting}. In a physical
classroom, instructors continuously read facial and vocal cues and adapt in
real time. That regulatory loop largely disappears in mass online settings.

Affective computing and educational psychology establish that boredom,
confusion, and frustration are not peripheral but integral to learning
dynamics~\cite{picard1997affective,pekrun2006control,dmello2012dynamics,baker2010better,dmello2015review}.
Automatic engagement recognition~\cite{whitehill2014faces,gupta2016daisee} and
speech emotion resources such as \textsc{IEMOCAP}~\cite{busso2008iemocap},
together with self-supervised speech encoders~\cite{chen2022wavlm} and
LLMs~\cite{brown2020language,kasneci2023chatgpt,lewis2020retrieval}, make an
end-to-end affect-aware tutor technically plausible.

\textbf{EmoLearn} implements that loop as a \emph{simulation research
prototype}: engagement signals come from annotated corpora
(\textsc{DAiSEE}, \textsc{IEMOCAP}), not live MOOC capture. The current system
contributes:
\begin{enumerate}
\item A product facial judge based on \textbf{FER+AU} (7 emotion scores + 10 AU
proxies $\rightarrow$ Random Forest), with a fine-tuned ViT used only as a
frame filter / fallback---not as the fusion input when AU is available.
\item A default \textbf{WavLM} speech stack (DAiSEE head + IEMOCAP ER hybrid),
replacing earlier wav2vec2 runtime defaults.
\item \textbf{Acoustic voice pairing} (state filter + 12-D cosine profile) with
an offline index over $10{,}039$ IEMOCAP utterances.
\item Late fusion with $60\%/40\%$ priors \emph{reweighted by confidence}, plus
clip and multi-clip \textbf{sequence} modes with temporal urgency rules.
\item LLM generation with RAG, twenty pedagogical templates, an affect-aware
conversational tutor, and instrumented form/similarity evaluation.
\end{enumerate}

Section~\ref{sec:related} reviews related work. Section~\ref{sec:arch}
presents architecture. Section~\ref{sec:methods} details methods.
Section~\ref{sec:impl} notes implementation choices. Section~\ref{sec:exp}
describes experiments. Section~\ref{sec:results} reports results.
Section~\ref{sec:discussion} discusses limits. Section~\ref{sec:conclusion}
concludes.

%==============================================================================
\section{Related Work}
\label{sec:related}
%==============================================================================

\subsection{Engagement and Affect in Learning}
Discrete affective states during technology-mediated learning have been
extensively discussed~\cite{dmello2015review,pekrun2006control,dmello2012dynamics,baker2010better}.
\textsc{DAiSEE}~\cite{gupta2016daisee} provides four engagement-related states
(Boredom, Confusion, Engagement, Frustration) with intensity levels, which we
adopt as the system's decision ontology.

\subsection{Facial and Vocal Sensing}
Facial engagement recognition~\cite{whitehill2014faces} and Facial Action
Coding System (FACS) inspired features motivate combining categorical emotion
scores with muscle-related cues. MediaPipe Face Landmarker blendshapes provide
practical AU \emph{proxies} (not human FACS coding). For speech, wav2vec
2.0~\cite{baevski2020wav2vec2} and WavLM~\cite{chen2022wavlm} offer strong
self-supervised backbones; EmoLearn's current default is WavLM with a frozen
backbone and a trained classification head for four \textsc{DAiSEE} states,
augmented by an IEMOCAP categorical ER anchor in a hybrid rule.

\subsection{Generative Adaptive Tutoring}
LLMs enable on-demand explanations and exercises~\cite{kasneci2023chatgpt};
RAG~\cite{lewis2020retrieval} grounds generation in a pedagogical store.
Design-science research~\cite{hevner2004design} motivates building measurable
artifacts. EmoLearn's claim is architectural: closing the loop from multimodal
perception to decision to generation, with dual pipelines (product vs.\ aligned
testbed) and an honest evaluation of pairing and mapping assumptions.

%==============================================================================
\section{System Architecture}
\label{sec:arch}
%==============================================================================

Figure~\ref{fig:pipeline} summarizes the product flow.

\begin{figure}[!t]
\centering
\fbox{\parbox{0.92\columnwidth}{\centering\small
\textbf{Scenario} (DAiSEE video clip + paired IEMOCAP audio)\\[3pt]
$\Big\downarrow$\\[3pt]
\textbf{Perception:} FER+AU (face) + WavLM SER (voice)\\
$\rightarrow$ fusion $60\%/40\%$ $\times$ confidence\\[3pt]
$\Big\downarrow$\\[3pt]
\textbf{Decision:} trigger? type A/B/C/D + template A1--D5\\
(+ temporal rules on sequences of 5 same-learner clips)\\[3pt]
$\Big\downarrow$\\[3pt]
\textbf{Generation:} LLM (GPT-OSS 20B via Groq) + RAG\\
+ affect-aware conversational tutor\\[3pt]
$\Big\downarrow$\\[3pt]
\textbf{Delivery:} React UI / PDF / SQLite journal
}}
\caption{Current EmoLearn product pipeline (simulation prototype).}
\label{fig:pipeline}
\end{figure}

\subsection{Two Pipelines, One Fusion Engine}
\begin{itemize}
\item \textbf{Product pipeline.} FastAPI + React. Face from real \textsc{DAiSEE}
video with FER+AU; voice from \textsc{IEMOCAP} audio \emph{paired} by DAiSEE
state filter and acoustic/affective profile similarity (not the learner's own
voice).
\item \textbf{Aligned testbed.} Offline validation where face and voice belong
to the \emph{same} IEMOCAP speaker and are synchronized. Both pipelines share
the same fusion engine.
\end{itemize}

\subsection{Runtime Decision Table}
Table~\ref{tab:runtime} states what actually decides in the deployed
configuration.

\begin{table}[!t]
\caption{Current product runtime judges.}
\label{tab:runtime}
\centering
\begin{tabular}{@{}p{0.22\columnwidth}p{0.68\columnwidth}@{}}
\toprule
\textbf{Brick} & \textbf{Production judge / encoder} \\
\midrule
Face $\rightarrow$ 4 states &
\textbf{FER+AU}: pretrained ViT (7 emotions) + 10 MediaPipe AU proxies
$\rightarrow$ Random Forest \\
Face (assistant) &
Fine-tuned ViT on \textsc{DAiSEE}: \textbf{frame filter} and reference
prediction; does \emph{not} choose the fused state when the AU classifier
is active \\
Voice $\rightarrow$ 4 states &
WavLM DAiSEE \textbf{head-only} fine-tune + hybrid IEMOCAP emotion-recognition
anchor \\
Voice pairing &
State filter + classical cosine on 12-D profile; acoustic index over
$10{,}039$ utterances \\
Fusion &
Late fusion $60\%$ face / $40\%$ voice $\times$ confidence \\
LLM &
Groq \textbf{GPT-OSS 20B} + ChromaDB RAG \\
\bottomrule
\end{tabular}
\end{table}

\subsection{Service Surface}
The prototype exposes single-clip analysis, multi-clip sequence analysis with
temporal rules, an affect-aware conversational tutor, scenario and template
catalogs, session history, IEMOCAP audio playback (corpus label vs.\ SER
preview), and health checks confirming WavLM as the speech backbone.

%==============================================================================
\section{Methods}
\label{sec:methods}
%==============================================================================

\subsection{Facial Perception: FER+AU (Product Judge)}
\label{sec:fer}

Faces are detected with MediaPipe \textbf{BlazeFace} (short/full range, with
mirroring). Five frames are sampled. Each frame yields:
\begin{itemize}
\item seven emotion probabilities from a pretrained Vision Transformer for
facial emotion recognition (seven categorical classes);
\item a 10-dimensional AU-proxy vector from MediaPipe Face Landmarker
blendshapes (brows, mouth, jaw, \ldots)---an approximation of FACS, not human
coding.
\end{itemize}
The concatenated $17$-D vector is classified by a Random Forest trained with
GridSearchCV on a stratified \textsc{DAiSEE} sample. AU features particularly
help separate \textbf{Frustration} and \textbf{Confusion}, which categorical
FER alone often confuses.

\paragraph{Fine-tuned ViT (not the product judge when AU is on).}
A ViT fine-tuned directly to four \textsc{DAiSEE} states is available. With AU
deployed, it filters low-confidence frames and provides a reference prediction;
it becomes the judge only if the AU-based classifier is unavailable. Best
measured full fine-tune reaches macro-F1 $0.388$ (Table~\ref{tab:vit}).

\subsection{Speech Perception: WavLM + Hybrid Anchor}
\label{sec:ser}

Default runtime uses a WavLM base encoder with a DAiSEE classification head
(backbone frozen, $800$ samples, augmentation) at accuracy $0.453$ /
macro-F1 $0.434$. A full unfreeze ablation ($1{,}200$ samples) reaches slightly
higher accuracy ($0.457$) but \textbf{worse} macro-F1 ($0.396$) and is optional
only. An IEMOCAP emotion-recognition head (neutral/happy/angry/sad, F1
$0.512$) acts as a hybrid anchor: the DAiSEE fine-tune is primary
($\approx 78\%$); the anchor is used especially when the fine-tune predicts
Confusion; if the fine-tune predicts Engagement while the anchor predicts
Frustration (e.g., laughter confused with anger), the fine-tune is kept.
Acoustic heuristics remain a complement and fallback.

\paragraph{Historical note on the fusion testbed.}
The aligned audiovisual bank ($n{=}80$) was measured when the learned SER
baseline was wav2vec2 ($0.487$ accuracy). Product runtime has since moved to
WavLM; we report both facts without conflating them.

\subsection{Acoustic Voice Pairing (Product Path)}
\label{sec:pair}

Because product face clips have no learner audio, IEMOCAP waveforms are paired
as follows: (1)~filter candidates by dominant DAiSEE state; (2)~rank by
classical cosine similarity on a 12-D profile (4 B/C/E/F scores, Baker/Russell
scenario VAD with neutral dominance $0.5$, duration, and 4 acoustic dims). An
offline index stores energy, pitch, ZCR, and brightness for all $10{,}039$
utterances with audio. Without the index, acoustic dims on candidates were
zero and pairing relied mainly on state+VAD. Scenario-side acoustic dims remain
a \emph{proxy} derived from engagement scores (DAiSEE clips have no audio)---a
refinement of pairing, not speaker synchrony.

\subsection{Fusion with Confidence Reweighting}
\label{sec:fusion}

Priors are $w_f{=}0.6$, $w_v{=}0.4$. Effective weights are
$w' \propto \mathrm{prior}\times\mathrm{confidence}$, then renormalized.
Pedagogical-priority fallback applies only when both confidences are null.
Hesitant FER frames are rejected before averaging (confidence threshold
$0.35$). On the aligned testbed, weighted vote and pedagogical priority remain
the compared strategies for accuracy vs.\ urgent recall.

\subsection{Temporal Rules and Sequence Mode}
Clip mode simulates an instant. Sequence mode simulates a mini-session: five
excerpts from the \emph{same} \textsc{DAiSEE} learner, each analyzed with the
full pipeline (independent IEMOCAP pairings per clip). Rules inspect the
history of \emph{detected} scores:

\begin{itemize}
\item \textbf{Immediate} (1 clip): e.g., Confusion $\ge 3$ and Frustration
$\ge 2$, or Engagement $=0$;
\item \textbf{Fast} (2--3 consecutive): Boredom/Confusion/Frustration $\ge 2$;
\item \textbf{Deferred} (4--5 clips): low Engagement or prolonged mild Boredom;
\item \textbf{None:} stable Engagement $\ge 2$ / acceptable zones.
\end{itemize}

\subsection{Pedagogical Decision and Generation}
Four content types map to twenty templates: \textbf{A} alternative explanation,
\textbf{B} concrete example, \textbf{C} progressive exercise, \textbf{D}
encouragement. On a 50-scenario benchmark, frequent templates include D1, C1,
D2/D4, A1, B1---coverage across A--D is effective.

Generation uses Groq \textbf{GPT-OSS 20B} (temperature $0.3$, retry/backoff)
with optional ChromaDB RAG over multilingual sentence embeddings. The
conversational tutor adapts tone, instruction, and temperature to the detected
state. Content quality is instrumented with form metrics and internal
BLEU/ROUGE similarity (template-median references---not expert gold).

\subsection{Technology Stack}
Python~3.11, FastAPI, SQLite, React; scikit-learn, PyTorch, Hugging Face
Transformers, OpenCV, \textbf{MediaPipe} (BlazeFace + Face Landmarker),
librosa; ChromaDB; pytest covering API, fusion, SER, temporal rules, RAG, and
training utilities.

\subsection{Illustrative Product Walkthrough}
Consider a \textsc{DAiSEE} clip labeled with rising Confusion. BlazeFace
extracts five frames; low-confidence frames are discarded by the fine-tuned ViT
filter. Surviving frames produce 7-D emotion scores and 10-D AU proxies; the
Random Forest outputs a four-state distribution dominated by Confusion. In
parallel, the pairing module filters IEMOCAP candidates to Confusion-mapped
emotions, ranks them by 12-D cosine similarity (now using real acoustic
features from the offline index), and feeds the selected waveform to WavLM.
Hybrid SER may consult the IEMOCAP ER anchor if Confusion is uncertain. Fusion
reweights the $60/40$ prior by confidence and yields a fused state. If temporal
history (sequence mode) shows two consecutive degraded clips, a \textbf{fast}
trigger fires; the decision engine selects, for example, template C1 (guided
exercise). RAG retrieves a short pedagogical snippet; GPT-OSS~20B generates the
text; the journal stores modalities, pairing method, template id, and
latencies. The UI can separately show the IEMOCAP corpus emotion label and the
WavLM prediction on the same audio---making annotation vs.\ inference explicit.

%==============================================================================
\section{Implementation Notes}
\label{sec:impl}
%==============================================================================

Perception modules expose probability distributions whenever possible so that
fusion operates on masses rather than brittle one-hot labels. Environment
flags select FER mapping mode, SER hybrid behavior, and whether the optional
full WavLM checkpoint is loaded. SQLite persists sessions, per-clip
predictions, fusion outcomes, selected templates, and generation metadata---a
requirement for auditability and for reproducing thesis tables.

LLM latency dominates end-to-end time; temperature $0.3$ stabilizes tone, and
exponential backoff absorbs Groq rate limits. Because the prototype is
simulation-oriented, reproducibility (fixed seeds, fixed aligned subset
construction, documented checkpoints) is prioritized over throughput.

Expert-evaluation HTML forms (students, teachers, experts) and offline
notebooks (OS1--OS4 exploratory analyses) accompany the software artifact but
are outside the runtime critical path.

%==============================================================================
\section{Experimental Protocol}
\label{sec:exp}
%==============================================================================

\paragraph{Corpora.}
\textsc{DAiSEE} is the engagement backbone. \textsc{IEMOCAP} supplies voice and
the aligned fusion subset. A pedagogical corpus of $\sim\!200$ generated items
supports form evaluation; $\sim\!60$ items are prepared for expert teacher
scoring (forms ready, scores pending).

\paragraph{IEMOCAP$\rightarrow$DAiSEE mapping (hypothesis).}
hap/exc$\rightarrow$Engagement; neu/sad$\rightarrow$Boredom;
ang/fru/dis$\rightarrow$Frustration; sur/fea$\rightarrow$Confusion. This is a
construct hypothesis, not a proven equivalence.

\paragraph{Evaluation layers.}
(1)~FER+AU and ViT ablations on \textsc{DAiSEE};
(2)~WavLM SER fine-tunes;
(3)~aligned fusion bank $n{=}80$ (Session~1 IEMOCAP, $\ge 1$\,s, 5 frames,
seed~42, $100\%$ faces detected);
(4)~product-proxy fusion on $n{=}50$ scenarios vs.\ dominant DAiSEE state;
(5)~acoustic pairing A/B on $50$ scenarios;
(6)~weight and mapping sensitivity;
(7)~LLM form ($N{=}208$) and BLEU/ROUGE ($N{=}189$);
(8)~synthetic reliability cases for confidence-weighted fusion ($n{=}6$).

Expert Delphi (L1) and teacher content ratings (L2) have instruments ready;
user piloting (L4, $n{\approx}30$) remains future work. Automatic metrics
(L3) are completed.

%==============================================================================
\section{Results}
\label{sec:results}
%==============================================================================

\subsection{Product Face Judge: FER+AU}
\textsc{DAiSEE} is imbalanced ($\sim\!75\%$ Engagement). Table~\ref{tab:ferau}
shows that FER+AU reaches macro-F1 $0.37$ with Frustration recall $34\%$,
versus a heuristic 7$\rightarrow$4 baseline at F1 $0.19$. Stratified RF on
emotion scores alone improves minority recall but collapses accuracy on a
balanced test; AU is retained as the production trade-off.

\begin{table}[!t]
\caption{Facial mapping to \textsc{DAiSEE} (product-relevant variants).}
\label{tab:ferau}
\centering
\begin{tabular}{@{}lccc@{}}
\toprule
\textbf{Variant} & \textbf{Acc.} & \textbf{F1} & \textbf{Frust.\ rec.} \\
\midrule
Imbalanced baseline collapse & $0.72$ & $0.26$ & $0\%$ \\
Oversampling only & $0.70$ & $0.30$ & $0\%$ \\
Stratified RF (7 scores) & $0.28^{*}$ & $0.27$ & $20\%$ \\
\textbf{FER+AU (runtime)} & $0.37$ & $\mathbf{0.37}$ & $\mathbf{34\%}$ \\
Heuristic 7$\rightarrow$4 & $0.26$ & $0.19$ & --- \\
\bottomrule
\end{tabular}
\vspace{2pt}
{\raggedright\footnotesize $^{*}$balanced test; Confusion recall $41\%$.\par}
\end{table}

\begin{table}[!t]
\caption{Fine-tuned ViT ablations (filter/fallback, not AU judge).}
\label{tab:vit}
\centering
\begin{tabular}{@{}lccc@{}}
\toprule
\textbf{Variant} & \textbf{Samples} & \textbf{Acc.} & \textbf{F1} \\
\midrule
Partial (2 layers) & $800$ & $0.36$ & $0.361$ \\
Large (4 layers) & ${\approx}1566$ & $0.378$ & $0.370$ \\
\textbf{Full FT (best)} & $800$ & $\mathbf{0.39}$ & $\mathbf{0.388}$ \\
\bottomrule
\end{tabular}
\end{table}

\subsection{WavLM Speech Fine-Tuning}
Table~\ref{tab:ser} compares SER checkpoints. The head-only DAiSEE model is the
default; full unfreeze is a negative ablation on macro-F1. The IEMOCAP ER
anchor supports the hybrid rule.

\begin{table}[!t]
\caption{WavLM fine-tunes (best epoch metrics).}
\label{tab:ser}
\centering
\begin{tabular}{@{}lccc@{}}
\toprule
\textbf{Model} & \textbf{$N$} & \textbf{Acc.} & \textbf{F1} \\
\midrule
\textbf{WavLM DAiSEE head (default)} & $800$ & $\mathbf{0.453}$ & $\mathbf{0.434}$ \\
WavLM DAiSEE full unfreeze & $1200$ & $0.457$ & $0.396$ \\
WavLM IEMOCAP ER anchor & $1200$ & $0.533$ & $0.512$ \\
\bottomrule
\end{tabular}
\end{table}

\subsection{Aligned Fusion Testbed ($n{=}80$)}
Table~\ref{tab:fusion} reports the speaker-aligned bank. Learned SER
(wav2vec2 at measurement time) nearly doubles heuristic voice accuracy
($+0.224$). Weighted-vote fusion reaches $0.500$ accuracy (above the best
unimodal $0.487$); pedagogical priority maximizes urgent recall ($0.400$).

\begin{table}[!t]
\caption{Aligned audiovisual fusion ($n{=}80$; SER learned $=$ wav2vec2 at testbed time).}
\label{tab:fusion}
\centering
\begin{tabular}{@{}lcc@{}}
\toprule
\textbf{Approach} & \textbf{Acc.} & \textbf{Urgent recall} \\
\midrule
Face --- FER 1 frame & $0.375$ & $0.250$ \\
Face --- FER multi-frame & $0.412$ & $0.275$ \\
Voice --- SER heuristic & $0.263$ & $0.225$ \\
Voice --- SER learned (wav2vec2) & $0.487$ & $0.200$ \\
Fusion --- pedagogical priority & $0.338$ & $\mathbf{0.400}$ \\
Fusion --- weighted vote & $\mathbf{0.500}$ & $0.125$ \\
\bottomrule
\end{tabular}
\end{table}

\subsection{Product-Proxy Pipeline ($n{=}50$)}
Against dominant \textsc{DAiSEE} state (not gold multimodal labels), product
fusion reaches accuracy $0.50$ and urgent recall $1.00$. Voice pairing matches
the target state in $50/50$ scenarios ($100\%$), versus a prior bias near
$58\%$ Engagement without the state filter.

\subsection{Acoustic Pairing A/B}
Adding real acoustic features raises mean profile similarity from $0.739$ to
$0.759$ ($\Delta{+}0.020$), changes the selected voice in $40/50$ scenarios
($80\%$), with $94\%$ of deltas positive, while state alignment stays at
$100\%$ (Table~\ref{tab:ab}).

\begin{table}[!t]
\caption{Acoustic pairing A/B ($50$ scenarios).}
\label{tab:ab}
\centering
\begin{tabular}{@{}lcc@{}}
\toprule
\textbf{Metric} & \textbf{VAD-only (A)} & \textbf{+ acoustic (B)} \\
\midrule
Mean similarity & $0.739$ & $\mathbf{0.759}$ \\
Voices changed & --- & $40/50$ ($80\%$) \\
Positive $\Delta$ & --- & $94\%$ \\
Target-state alignment & $100\%$ & $100\%$ \\
\bottomrule
\end{tabular}
\end{table}

\subsection{Weight and Mapping Sensitivity}
On the aligned set, face/voice weights were swept from $0.5/0.5$ to $0.9/0.1$
(Table~\ref{tab:weights}). The default $0.6/0.4$ maximizes urgent recall
($0.425$) and is retained; $0.7/0.3$ yields slightly higher accuracy
($0.338$). Remapping IEMOCAP neutral$\rightarrow$Engagement changes $17.0\%$
of states across $10{,}039$ utterances; sadness$\rightarrow$Frustration
changes $10.8\%$; surprise and disgust remappings are minor or null
(Table~\ref{tab:map}).

\begin{table}[!t]
\caption{Fusion weight sensitivity (aligned $n{=}80$).}
\label{tab:weights}
\centering
\begin{tabular}{@{}lcc@{}}
\toprule
\textbf{Face / voice} & \textbf{Acc.} & \textbf{Urgent recall} \\
\midrule
$0.5$ / $0.5$ & $0.287$ & $0.375$ \\
$\mathbf{0.6}$ / $\mathbf{0.4}$ (default) & $0.300$ & $\mathbf{0.425}$ \\
$0.7$ / $0.3$ & $\mathbf{0.338}$ & $0.400$ \\
$0.8$ / $0.2$ & $0.325$ & $0.350$ \\
$0.9$ / $0.1$ & $0.325$ & $0.350$ \\
\bottomrule
\end{tabular}
\end{table}

\begin{table}[!t]
\caption{IEMOCAP$\rightarrow$DAiSEE mapping sensitivity ($10{,}039$ utterances).}
\label{tab:map}
\centering
\begin{tabular}{@{}lc@{}}
\toprule
\textbf{Alternative mapping} & \textbf{\% states changed} \\
\midrule
neutral $\rightarrow$ Engagement (vs.\ Boredom) & $17.0\%$ \\
sadness $\rightarrow$ Frustration (vs.\ Boredom) & $10.8\%$ \\
surprise $\rightarrow$ Engagement (vs.\ Confusion) & $1.1\%$ \\
disgust $\rightarrow$ Boredom (vs.\ Frustration) & $0.0\%$ \\
\bottomrule
\end{tabular}
\end{table}

\subsection{LLM Form and Internal Similarity}
On $N{=}208$ generations: mean length $222$ words (min $113$, max $410$),
Kandel--Moles readability $95.0$, mean repetition rate $0.479$, length-target
hit rate $0.832$. Internal similarity to per-template median references
($N{=}189$): BLEU-4 $0.139$, ROUGE-1 F1 $0.449$, ROUGE-2 $0.176$, ROUGE-L
$0.272$. These measure internal coherence, \emph{not} pedagogical
correctness.

\subsection{Confidence-Weighted Fusion Reliability}
On six synthetic critical cases (noisy modality, total uncertainty, \ldots),
continuous weighting $w=\mathrm{prior}\times\mathrm{confidence}$ achieves
accuracy $1.0$ on those crafted cases---e.g., down-weighting a noisy voice
instead of forcing a wrong urgent state via pedagogical priority. This is a
reliability unit test, not a claim of perfect real-world accuracy.

\subsection{Measured Gains (Aligned Bank)}
Relative to the aligned $n{=}80$ table: learned SER vs.\ heuristic
$+0.224$ accuracy; multi-frame FER vs.\ one frame $+0.037$; weighted-vote
fusion vs.\ pedagogical priority $+0.162$ accuracy; weighted vote vs.\ best
unimodal $+0.013$. The last gain is small but positive---fusion helps under
aligned pairing, without miracle margins.

%==============================================================================
\section{Discussion and Limitations}
\label{sec:discussion}
%==============================================================================

EmoLearn is scoped as a simulation prototype. Key tensions:

\begin{itemize}
\item \textbf{Mapping} emotions$\rightarrow$engagement remains a hypothesis;
sensitivity shows non-trivial impact for neutral/sadness.
\item \textbf{Product pairing} is cross-corpus (DAiSEE face + IEMOCAP voice),
mitigated scientifically by the aligned testbed and operationally by state
filtering + acoustic profiles---still not the learner's own voice.
\item \textbf{Scenario acoustic dims} are proxies; A/B evaluates profile
similarity, not perceptual ``best voice'' judgments.
\item \textbf{Templates} are an argued design; expert Delphi and teacher
ratings instruments are ready but scores are not yet collected.
\item \textbf{LLM substance} needs human validation; automatic metrics cover
form and internal similarity only.
\item \textbf{Testbed SER} figures use wav2vec2-era learned speech; runtime is
WavLM---re-running the aligned bank with WavLM is required for strict
comparability.
\item \textbf{Breadth vs.\ depth:} the claimed contribution is end-to-end
architecture with honest limits, not a single SOTA unimodal score.
\end{itemize}

Table~\ref{tab:limits} consolidates status.

\begin{table}[!t]
\caption{Assumed tensions and status.}
\label{tab:limits}
\centering
\begin{tabular}{@{}p{0.30\columnwidth}p{0.60\columnwidth}@{}}
\toprule
\textbf{Tension} & \textbf{Status} \\
\midrule
Fusion vs.\ unimodal & Measured on aligned bank + confidence weighting \\
Emotion$\rightarrow$engagement map & Hypothesis + sensitivity \\
Adaptive vs.\ simulation & Prototype (no live MOOC capture) \\
Heterogeneous face/voice pairing & Documented; acoustic index + aligned subset \\
20 templates & Design argued; expert forms pending scores \\
LLM content quality & Form/similarity done; substance pending humans \\
SER testbed vs.\ runtime & wav2vec2 bank vs.\ WavLM product --- to re-align \\
\bottomrule
\end{tabular}
\end{table}

Ethically, any live deployment would require consent, minimization, and an
opt-out from sensing while retaining non-affective tutoring. Decision journals
already support inspectability of triggers, modalities, pairing method, and
templates. Instructors should be able to understand, at least at a high level,
why an adaptation was produced.

%==============================================================================
\section{Conclusion}
\label{sec:conclusion}
%==============================================================================

We presented the current EmoLearn prototype: FER+AU facial judging, WavLM
hybrid speech recognition, confidence-reweighted $60/40$ fusion, acoustic
IEMOCAP pairing, temporal sequence rules, and GPT-OSS~20B generation with RAG
and an affect-aware tutor. Empirically, FER+AU improves macro-F1 over
heuristics ($0.37$ vs.\ $0.19$); WavLM head-only is preferred to full unfreeze
on F1 ($0.434$ vs.\ $0.396$); weighted-vote fusion reaches $0.500$ on the
aligned bank; acoustic pairing improves mean similarity ($+0.020$) while
preserving $100\%$ state alignment; and automatic content metrics characterize
form without claiming pedagogical gold quality.

Future work includes collecting L1/L2 expert scores, a user pilot
($n{\approx}30$), live synchronized capture under ethics approval, re-evaluating
the aligned fusion bank with WavLM, fairness analyses across recording
conditions, and outcome-oriented studies (comprehension, persistence) beyond
classification accuracy.

\section*{Acknowledgment}
The authors thank the thesis supervisor and the providers of the \textsc{DAiSEE}
and \textsc{IEMOCAP} corpora under academic agreements.

\bibliographystyle{IEEEtranN}
\bibliography{references}
\balance

\end{document}
